\documentclass[11pt]{book}
\usepackage{fontspec}
\setmainfont{Latin Modern Roman}
\setsansfont{Latin Modern Sans}
\usepackage{amsmath,amssymb,mathtools}
\usepackage{unicode-math}
\setmathfont{Latin Modern Math}
\usepackage{xurl}
\usepackage[hidelinks]{hyperref}
\newcommand{\sref}[2]{\hyperlink{source:#1:#2}{[S#2]}}
\newcommand{\cataloguescope}[1]{\paragraph{Mathematical question.}#1}
\begin{document}
\section{Research Problems in Function Theory — Problem 3.26}\label{statement}

\subsection{Definitions and mathematical statement}

\paragraph{Shared notation.}$\mathbb N=\{1,2,\ldots\}$, and $\mathbb Z,\mathbb Q,\mathbb R,\mathbb C$ denote the integers, rationals, reals and complex numbers. Any other conventions are specified locally.


For $n\ge2$, let $B_n=\{x\in\mathbb R^n:|x|<1\}$. A subharmonic function is an upper semicontinuous function satisfying the sub-mean inequality on every ball compactly contained in its domain. Here require $u:B_n\to[0,1)$ and $u(0)>0$. An admissible path is a rectifiable continuous curve $\gamma:[0,1]\to\overline B_n$ with $\gamma(0)=0$, $\gamma(1)\in\partial B_n$, $\gamma(t)\in B_n$ and $u(\gamma(t))>0$ for $t<1$. Its Euclidean arc length is denoted by $\ell(\gamma)$. No boundary value of $u$ is required.
\paragraph{Statement recorded in the supplied source.}
For each fixed $n\ge4$, find a continuous function $f_n:(0,1)\to(0,\infty)$ such that every such $u$ has an admissible path with
\[\ell(\gamma)\le f_n(u(0)).\]
For $n=2$, determine the smallest exponent $\beta$ for which one can use a uniform bound $f_2(t)=Ct^{-\beta}$, $0<t<1$, with a constant $C$ independent of $u,t$; determine the optimal infimum of the exponents and whether it is attained. \sref{statement}{2}
\subsection{Short English statement}
Bound the length of a positive path from the centre to the boundary for bounded subharmonic functions in dimensions at least four, and find the optimal power exponent for the planar case.
\subsection{Sources}
\begin{description}
\item[\hypertarget{source:statement:1}{S1}] ULAM.AI and the UnsolvedMath Contributors, \emph{UnsolvedMath: A Curated Collection of Open Mathematics Problems}, catalogue entry. CC BY 4.0. \url{https://huggingface.co/datasets/ulamai/UnsolvedMath}.
\item[\hypertarget{source:statement:2}{S2}] Walter K. Hayman and Edward F. Lingham, \emph{Research Problems in Function Theory}, arXiv:1809.07200, 2018. Problem 3.26, its complete Update 3.26, and the relevant chapter notation. \url{https://arxiv.org/pdf/1809.07200}.
\end{description}
\label{end:statement}
\end{document}
